\chapter{Dialysis Catheter Placement}

\section*{Introduction \& Clinical Indications}
Dialysis catheter placement provides temporary or longer-term vascular access for hemodialysis or continuous renal replacement therapy. A non-tunneled catheter is used when access is needed urgently or for a short period, while a tunneled cuffed catheter is used when a fistula or graft cannot yet be created or used. Catheters are generally a bridge or alternative to an arteriovenous fistula or graft because infection and thrombosis risks are higher.

Indications include severe kidney failure requiring urgent dialysis, delayed maturation of permanent access, and selected patients in whom permanent access is unsuitable. In an emergency, life-threatening hyperkalemia, acidosis, pulmonary edema, or uremic complications take priority over elective access planning. Catheter choice should consider expected duration, dialysis prescription, previous access, and local expertise.

\section*{Relevant Surgical Anatomy}
The right internal jugular vein is commonly preferred because it offers a relatively straight route to the superior vena cava and right atrium. The left internal jugular, femoral veins, and, less often, subclavian vein may be alternatives. Subclavian access is usually avoided when possible because central venous stenosis can compromise future fistula or graft options. The carotid artery, pleura, thoracic duct on the left, and mediastinal vessels are important structures at risk.

\section*{Preoperative Workup \& Preparation}
Assess urgency, platelet count, coagulation status, infection, allergies, prior central lines, and the condition of potential access sites. Ultrasound guidance is standard for internal-jugular cannulation when available. Review imaging when central venous stenosis or thrombosis is suspected. Explain bleeding, arterial puncture, pneumothorax, malposition, infection, thrombosis, arrhythmia, and catheter dysfunction. Use maximal sterile barriers, skin antisepsis, and a site-specific plan for antibiotics according to institutional policy.

\section*{Surgical Technique \& Procedure}
With the patient monitored and positioned safely, the operator identifies the target vein with ultrasound and confirms compressibility and anatomy. Venous puncture is followed by guidewire passage, tract dilation, and catheter advancement using the Seldinger technique. The catheter tip is positioned in an appropriate central location, lumens are aspirated and flushed, and the device is secured with a sterile dressing. A tunneled catheter is passed through a subcutaneous tract and its cuff is seated in the tunnel to encourage tissue ingrowth.

The guidewire must remain controlled throughout dilation and catheter exchange. Resistance, abnormal wire position, ectopy, or poor aspiration warrants stopping and reassessment. Tip confirmation and post-procedure imaging are performed according to catheter type, insertion site, and local protocol; ultrasound can assess complications at the neck or groin. The catheter should not be used until position and function are acceptable.

\section*{Complications \& Risks}
Early complications include arterial puncture, hematoma, pneumothorax, hemothorax, air embolism, arrhythmia, malposition, and guidewire loss. Later problems include catheter-related bloodstream infection, tunnel infection, fibrin sheath, thrombosis, central venous stenosis, and inadequate blood flow. Fever, rigors during dialysis, exit-site drainage, chest pain, dyspnea, or sudden loss of function require prompt assessment.

\section*{Postoperative Care \& Recovery}
Inspect the insertion and exit sites, document the catheter type and tip position, and teach hand hygiene, dressing care, and warning signs. Avoid unnecessary manipulation and maintain a closed system. Dialysis staff should follow an access-preservation protocol, including appropriate locking solution and prompt culture-based treatment when infection is suspected. A permanent access plan should be revisited early rather than allowing an avoidable long-term catheter dependence.

\section*{Outcomes \& Prognosis}
Placement is usually successful when ultrasound guidance, sterile technique, and appropriate imaging are used. Long-term outcome depends more on infection prevention, catheter patency, and transition to permanent access than on the insertion alone. The access plan should be individualized with nephrology, vascular access, and interventional teams.

\section*{References}
\begin{enumerate}
\item National Kidney Foundation. KDOQI Clinical Practice Guideline for Vascular Access.
\item Oliver MJ, et al. Canadian Society of Nephrology Clinical Practice Guidelines for Vascular Access.
\item CDC. Guidelines for the Prevention of Intravascular Catheter-Related Infections.
\end{enumerate}
